The Virgo cluster of galaxies is the nearest large cluster which extends over
nearly 10 degrees across the sky and contains a number of bright galaxies. It
is very interesting to find the distance to Virgo and to deduce certain
cosmological information from it. The table below provides the distance
estimates using various distance indicators (listed in the left column). The
right column lists the mean distance $d_i \pm$ the standard deviation $s_i$.

\begin{center}
\begin{tabular}{cl c}
$i$ & Distance Indicator & Virgo Distance (Mpc) \\
\hline
1 & Cepheids                       & $14.9 \pm 1.2$ \\
2 & Novae                          & $21.1 \pm 3.9$ \\
3 & Planetary Nebulae              & $15.2 \pm 1.1$ \\
4 & Globular Cluster               & $18.8 \pm 3.8$ \\
5 & Surface Brightness Fluctuation & $15.9 \pm 0.9$ \\
6 & Tully-Fisher relation          & $15.8 \pm 1.5$ \\
7 & Faber-Jackson relation         & $16.8 \pm 2.4$ \\
8 & Type Ia Supernovae             & $19.4 \pm 5.0$ \\
\end{tabular}
\end{center}

\begin{description}
\item[1.] By applying a weighted mean, compute the average distance (which can
  be taken as an estimate to the distance to Virgo)
  \[ d_{\mathrm{avg}}
     = \frac{\sum_i \dfrac{d_i}{s_i^2}}{\sum_i \dfrac{1}{s_i^2}} \]
  where the sum runs over the eight distance indicators used.
\item[2.] What is the uncertainty (rms) (in units of Mpc) in that estimate?
\item[3.] Spectra of the galaxies in Virgo indicate an average recession
  velocity of 1136\,km/sec for the cluster. Can you estimate the Hubble
  constant $H_0$ and its uncertainty (rms)?
\item[4.] What is the Hubble Time (age of the universe) using the value of the
  Hubble constant you found and the uncertainty (rms)?
\end{description}
